\section{Introduction : deux chemins, un m\^eme r\'eseau}

\subsection{L'observation}

\`A la fin du chantier multic\oe{}ur, la recherche UCI de LapZero-Chess
tournait entre 500 et 850 simulations par seconde selon la position et le
nombre de workers. Le r\'eseau \'etait un ResNet de 10 blocs et 128 filtres,
export\'e en ONNX, ex\'ecut\'e sur un GPU RTX PRO 2000 Blackwell portable avec
le provider CUDA d'ONNX Runtime. Le banc self-play, lui, mesurait 3,05\,ms pour
un lot de 8 positions, soit 0,38\,ms par position et environ 2\,600
\'evaluations par seconde, avec exactement le m\^eme mod\`ele et le m\^eme GPU.

Le tableau~\ref{tab:observation} r\'esume l'\'ecart : par position, la
recherche payait 1,4 \`a 2,9\,ms l\`a o\`u le self-play payait 0,38\,ms, un
facteur quatre \`a huit.

\begin{table}[h]
\centering
\small
\caption{Co\^ut par position avant le chantier. Recherche : phase
\'evaluateur de 700 simulations, chemin \code{step\_analysis}, lot 8, workers
8. Self-play : banc partag\'e, 8 plateaux, 100 vagues.}
\label{tab:observation}
\begin{tabular}{lrrr}
\toprule
Chemin & ms par position & sims/s & Source \\
\midrule
Self-play, lot 8 fixe & 0,38 & $\approx 2\,600$ & \code{selfplay-after-gpu.json} \\
Recherche, ouverture & 1,19 & 847 & \code{after.json} \\
Recherche, milieu & 2,87 & 349 & \code{after.json} \\
Recherche, finale & 1,77 & 565 & \code{after.json} \\
\bottomrule
\end{tabular}
\end{table}

\subsection{La question}

L'\'evaluateur est le m\^eme objet \code{ONNXEvaluator} dans les deux cas. La
diff\'erence ne pouvait donc pas venir du r\'eseau lui-m\^eme. Trois familles
d'hypoth\`eses restaient :

\begin{enumerate}[leftmargin=1.4em,itemsep=3pt]
  \item \textbf{La plomberie de l'appel.} La recherche reconstruit un vecteur
  d'entr\'ee par vague, redimensionne les sorties et appelle le m\^eme
  \code{evaluate\_batch}, mais avec des tailles de lot variables et des
  buffers diff\'erents \`a chaque appel. Le self-play appelle en boucle serr\'ee
  avec un buffer fixe.
  \item \textbf{Les allocations.} Chaque vague alloue un vecteur d'entr\'ee
  d'environ 200\,kio et laisse \code{evaluate\_batch} redimensionner les
  sorties. Le self-play r\'eutilise tout.
  \item \textbf{L'occupation GPU.} Le r\'eseau est minuscule : cartes 8$\times$8,
  128 filtres, environ 0,4\,GFLOP par position. Le GPU pouvait \^etre sous-occup\'e,
  et le co\^ut r\'eel pouvait venir des noyaux et non des donn\'ees.
\end{enumerate}

Les trois hypoth\`eses ne s'excluent pas, et l'\'etude de co\^ut
\chemin{2026-09-19-cout-calcul-cpu-gpu.md} les avait list\'ees sans pouvoir
trancher, faute de mesure qui s\'epare l'appel de l'arbre.

\subsection{La m\'ethode et les contributions}

La d\'emarche retenue est exp\'erimentale et hi\'erarchis\'ee : mesurer d'abord
l'\'evaluateur \emph{hors arbre}, puis ne corriger que la cause d\'emontr\'ee,
et refuser par la mesure toute piste qui ne passe pas la barri\`ere qualit\'e.

\begin{resume}[Les contributions de ce chantier]
\begin{itemize}[leftmargin=1.2em,itemsep=2pt]
  \item un banc C++ \code{evaluator\_batch\_bench} qui mesure la courbe du lot
  hors arbre, GPU et CPU, avec un mode \og lots altern\'es \fg{} qui reproduit
  la recherche, et des chronom\`etres internes s\'eparant \code{session->Run}
  du softmax ;
  \item la d\'ecouverte et la quantification de la re-planification d'ONNX
  Runtime au changement de forme : facteur quatre sur l'appel ;
  \item le lot de forme fixe, appliqu\'e aux vagues multic\oe{}ur et au chemin
  mono, s\'emantiquement neutre et activ\'e dans le bot UCI ;
  \item l'\'evaluation de deux pistes concurrentes, FP16 et r\'eglages de
  divergence, toutes deux refus\'ees sur preuves ;
  \item la mesure du cas chaud, qui corrige une r\'eserve des rapports
  pr\'ec\'edents sur la loi d'Amdahl ;
  \item un verdict chiffr\'e sur la production continue et le plafond restant.
\end{itemize}
\end{resume}

Le reste du rapport suit l'ordre : contexte et hypoth\`eses (chapitre~2),
conception et d\'ecouverte (chapitre~3), impl\'ementation (chapitre~4),
validation (chapitre~5), mesures (chapitre~6), discussion (chapitre~7),
conclusion (chapitre~8).
